\documentclass[11pt]{article}
\usepackage[utf8]{inputenc}
\usepackage[T1]{fontenc}
\usepackage{amsmath}
\usepackage{amssymb}
\usepackage{amsthm}
\usepackage[margin=1.15in]{geometry}
\usepackage{hyperref}

\newtheorem{theorem}{Theorem}
\newtheorem{lemma}{Lemma}
\theoremstyle{remark}
\newtheorem{remark}{Remark}

\title{Refutation of conjecture 00000000416\\[2mm]
\large The orbit count of evacuation on tableaux of shape $(n,n)$ is inconsistent at $n=2$}
\author{Submitted to The Last Math Competition}
\date{\today}

\begin{document}
\maketitle

\begin{abstract}
\noindent
Conjecture 00000000416 asserts that, for the evacuation operator acting on the
standard Young tableaux of two-row shape $(n,n)$, the number of orbits is
$2^{n-1}$, every orbit length divides $2n$, and the number of orbits of length
exactly $2$ is the Fibonacci number $F_{n-1}$. We show that the first and third
assertions are jointly impossible already at $n=2$: they require at least three
tableaux, whereas the shape $(2,2)$ carries exactly two. The argument uses
nothing about evacuation beyond the fact that an orbit is nonempty; in
particular it does not use that evacuation is an involution.
\end{abstract}

\section{The conjecture}

The statement under consideration, verbatim from the corpus, is the following.

\begin{quote}
\emph{Definition: The evacuation operator $e$ acts on standard Young tableaux of
two-row shape $(n,n)$; its orbits are the cycles of the iterates of $e$.
Conjecture: The number of orbits is $2^{n-1}$, every orbit length divides $2n$,
and the number of orbits of length exactly $2$ is the Fibonacci number
$F_{n-1}$.}
\end{quote}

Write $T_n$ for the set of standard Young tableaux of shape $(n,n)$, and let
\[
  C_n \;=\; \frac{1}{n+1}\binom{2n}{n}
\]
denote the $n$-th Catalan number. It is classical that $|T_n| = C_n$: a standard
tableau of shape $(n,n)$ is determined by the set of entries placed in the top
row, and the ballot condition says exactly that this set encodes a Dyck path of
semilength~$n$.

The conjecture makes three simultaneous claims about the action of $e$ on
$T_n$:
\begin{align}
  \#\{\text{orbits of } e\} &= 2^{n-1}, \label{eq:orbits}\\
  \text{every orbit length} &\mid 2n, \label{eq:divides}\\
  \#\{\text{orbits of length exactly } 2\} &= F_{n-1}. \label{eq:lengthtwo}
\end{align}

\section{The contradiction at $n=2$}

\begin{lemma}\label{lem:shape22}
The shape $(2,2)$ carries exactly two standard Young tableaux.
\end{lemma}

\begin{proof}
The entries are $1,2,3,4$. The top-left cell is smaller than both its row
neighbour and its column neighbour, and the bottom-right cell is larger than
both its row neighbour and its column neighbour; since the top-left entry is
below the bottom-right entry along the top row, the top-left entry is smaller
than all three others, hence equals~$1$, and symmetrically the bottom-right
entry equals~$4$. The remaining entries $2$ and $3$ fill the two off-diagonal
cells, and each of the two placements is admissible. Explicitly,
\[
  \begin{array}{cc} 1 & 2 \\ 3 & 4 \end{array}
  \qquad\text{and}\qquad
  \begin{array}{cc} 1 & 3 \\ 2 & 4 \end{array}
\]
are the only possibilities. (Equivalently $|T_2| = C_2 = 2$.)
\end{proof}

\begin{theorem}\label{thm:main}
Conjecture 00000000416 is false.
\end{theorem}

\begin{proof}
Suppose the conjecture holds at $n=2$. By~\eqref{eq:orbits} there are
$2^{2-1} = 2$ orbits in total. By~\eqref{eq:lengthtwo}, the number of orbits of
length exactly $2$ is $F_{2-1} = F_1 = 1$.

The orbits partition $T_2$. One orbit therefore contains exactly $2$ tableaux,
and the second orbit is nonempty, so it contains at least $1$. Hence
\[
  |T_2| \;\ge\; 2 + 1 \;=\; 3.
\]
But $|T_2| = 2$ by Lemma~\ref{lem:shape22}. This is a contradiction.

Therefore the conjunction of \eqref{eq:orbits} and \eqref{eq:lengthtwo} fails
at $n=2$, and with it the conjecture.
\end{proof}

\section{What the argument does and does not use}

The proof of Theorem~\ref{thm:main} uses exactly one property of evacuation:
that its orbits are nonempty. It does \emph{not} use that $e$ is an involution,
nor that orbit lengths are $1$ or $2$, nor any information about how $e$ acts on
a particular tableau. This is deliberate --- it makes the refutation robust
against any disagreement about the precise definition of evacuation, since any
map whatsoever whose iterates partition $T_2$ into two orbits, one of size $2$,
would require $|T_2| \ge 3$.

\begin{remark}[A second, independent failure]
If one additionally uses that evacuation is an involution, so that every orbit
has length $1$ or $2$, then \eqref{eq:orbits} and \eqref{eq:lengthtwo} force
\[
  |T_n| \;=\; \bigl(2^{n-1} - F_{n-1}\bigr)\cdot 1 \;+\; F_{n-1}\cdot 2
  \;=\; 2^{n-1} + F_{n-1}.
\]
For $n = 4$ the left side is $C_4 = 14$ while the right side is $8 + 2 = 10$, so
the conjecture fails again, independently of the $n = 2$ case. (For $n = 3$ the
identity $5 = 4 + 1$ does hold, and for $n = 1$ the identity $1 = 1 + 0$ holds
as well; the failures begin at $n = 2$ and resume from $n = 4$.)
\end{remark}

\begin{remark}[On claim \eqref{eq:divides}]
We make no assertion about \eqref{eq:divides}. It is consistent with the
others: orbit lengths are at most $2$ when $e$ is an involution, and both $1$
and $2$ divide $2n$ for $n \ge 1$. The refutation rests entirely on the
joint inconsistency of \eqref{eq:orbits} and \eqref{eq:lengthtwo}.
\end{remark}

\section{Formalisation}

The accompanying Lean~4 project formalises the argument in
\texttt{lean4/Main.lean}. It contains:

\begin{itemize}
  \item an exhaustive enumeration of the $4^4 = 256$ fillings of the diagram
        $(2,2)$ by the letters $0,1,2,3$, filtered to the standard ones. The
        theorem \texttt{sols\_length} states that exactly two survive, and is
        proved by \texttt{rfl}: the kernel performs the search and checks the
        count, so the theorem depends on \emph{no} axioms at all;
  \item the counting lemma \texttt{orbit\_bound}: if a list is split into two
        nonempty parts and one part has at least two elements, the total is at
        least~$3$;
  \item \texttt{refutation\_at\_two}: no such splitting of the tableaux of shape
        $(2,2)$ exists.
\end{itemize}

The project uses only the Lean~4 core (no Mathlib) and contains no
\texttt{sorry}. The axioms used are \texttt{propext} and \texttt{Quot.sound},
inherited from the standard treatment of lists.

\section{Reproducing the enumeration}

The accompanying \texttt{reproduce.py} (standard library only) enumerates the
standard Young tableaux of shape $(n,n)$ for small $n$, reports $C_n$ alongside
the quantity $2^{n-1} + F_{n-1}$ required by the conjecture, and prints the
tableaux of shape $(2,2)$. Running it reproduces the numbers used above.

\end{document}
